% GENERATED FILE. Edit canonical lessons, then run npm run generate:books.
% canonical-source-hash: fnv1a64:f26b97f3a499663f
% canonical-lessons: RU-C257-kompaniya, RU-C257-firma, RU-C257-proyekt, RU-C257-zadacha, RU-C257-opyt, RU-R257-first-pass-words-from-chapters-250-253, RU-R257-second-pass-words-from-chapters-254-257

\chapter{Company, Firm, Project, Task, Experience}
\label{ch:ru-company-firm-project-task-experience}
\hlchaptermodality{257}

\begin{chapteropening}
\textbf{By the end of this chapter:} \emph{I can say a company, a firm, a project, a task and experience.}

Complete the last lesson of chapter 257: I can say a company, a firm, a project, a task and experience.
\end{chapteropening}

\section[\texorpdfstring{\ru{компания} (kompániya)}{компания (kompániya)}]{\ru{компания} (kompániya) — a company}

\label{lesson:RU-C257-kompaniya}



\begin{quote}
Before the new one: say the Russian for a musician, then the Russian for an actor.
\end{quote}

\subsection*{You'll want to know: \ru{компания}}
\textbf{\ru{компания}} (\emph{kompániya}) — \textquotedblleft{}a company\textquotedblright{}.

\begin{grammarlens}[title={What you've built}]
One more word for this chapter.
\end{grammarlens}

\subsection*{Your turn}
\begin{itemize}
\raggedright
  \item \emph{Say it:} \emph{kompániya}
  \item \emph{Say it:} \emph{kompániya}, once more
  \item \emph{Say it:} say \emph{kompániya}
  \item \emph{Recall:} say the Russian for a musician, then the Russian for an actor, then say \emph{kompániya} again
\end{itemize}

\begin{tcolorbox}[breakable,colback=teal!4,colframe=teal!35!black,title={Before you move on}]
What does \ru{компания} mean? (\textquotedblleft{}A company\textquotedblright{}.) Say it once more.
\end{tcolorbox}
\section[\texorpdfstring{\ru{фирма} (fírma)}{фирма (fírma)}]{\ru{фирма} (fírma) — a firm}

\label{lesson:RU-C257-firma}



\begin{quote}
Before the new one: say the Russian for an actor, then the Russian for a company.
\end{quote}

\subsection*{You'll want to know: \ru{фирма}}
\textbf{\ru{фирма}} (\emph{fírma}) — \textquotedblleft{}a firm\textquotedblright{}.

\begin{grammarlens}[title={What you've built}]
One more word for this chapter.
\end{grammarlens}

\subsection*{Your turn}
\begin{itemize}
\raggedright
  \item \emph{Say it:} \emph{fírma}
  \item \emph{Say it:} \emph{fírma}, once more
  \item \emph{Say it:} say \emph{fírma}
  \item \emph{Recall:} say the Russian for an actor, then the Russian for a company, then say \emph{fírma} again
\end{itemize}

\begin{tcolorbox}[breakable,colback=teal!4,colframe=teal!35!black,title={Before you move on}]
What does \ru{фирма} mean? (\textquotedblleft{}A firm\textquotedblright{}.) Say it once more.
\end{tcolorbox}
\section[\texorpdfstring{\ru{проект} (proyékt)}{проект (proyékt)}]{\ru{проект} (proyékt) — a project}

\label{lesson:RU-C257-proyekt}



\begin{quote}
Before the new one: say the Russian for a company, then the Russian for a firm.
\end{quote}

\subsection*{You'll want to know: \ru{проект}}
\textbf{\ru{проект}} (\emph{proyékt}) — \textquotedblleft{}a project\textquotedblright{}.

\begin{grammarlens}[title={What you've built}]
One more word for this chapter.
\end{grammarlens}

\subsection*{Your turn}
\begin{itemize}
\raggedright
  \item \emph{Say it:} \emph{proyékt}
  \item \emph{Say it:} \emph{proyékt}, once more
  \item \emph{Say it:} say \emph{proyékt}
  \item \emph{Recall:} say the Russian for a company, then the Russian for a firm, then say \emph{proyékt} again
\end{itemize}

\begin{tcolorbox}[breakable,colback=teal!4,colframe=teal!35!black,title={Before you move on}]
What does \ru{проект} mean? (\textquotedblleft{}A project\textquotedblright{}.) Say it once more.
\end{tcolorbox}
\section[\texorpdfstring{\ru{задача} (zadácha)}{задача (zadácha)}]{\ru{задача} (zadácha) — a task; a problem (to solve)}

\label{lesson:RU-C257-zadacha}



\begin{quote}
Before the new one: say the Russian for a firm, then the Russian for a project.
\end{quote}

\subsection*{You'll want to know: \ru{задача}}
\textbf{\ru{задача}} (\emph{zadácha}) — \textquotedblleft{}a task; a problem (to solve)\textquotedblright{}.

\begin{grammarlens}[title={What you've built}]
One more word for this chapter.
\end{grammarlens}

\subsection*{Your turn}
\begin{itemize}
\raggedright
  \item \emph{Say it:} \emph{zadácha}
  \item \emph{Say it:} \emph{zadácha}, once more
  \item \emph{Say it:} say \emph{zadácha}
  \item \emph{Recall:} say the Russian for a firm, then the Russian for a project, then say \emph{zadácha} again
\end{itemize}

\begin{tcolorbox}[breakable,colback=teal!4,colframe=teal!35!black,title={Before you move on}]
What does \ru{задача} mean? (\textquotedblleft{}A task; a problem (to solve)\textquotedblright{}.) Say it once more.
\end{tcolorbox}
\section[\texorpdfstring{\ru{опыт} (ópyt)}{опыт (ópyt)}]{\ru{опыт} (ópyt) — experience; an experiment}

\label{lesson:RU-C257-opyt}



\begin{quote}
Before the new one: say the Russian for a project, then the Russian for a task; a problem.
\end{quote}

\subsection*{You'll want to know: \ru{опыт}}
\textbf{\ru{опыт}} (\emph{ópyt}) — \textquotedblleft{}experience; an experiment\textquotedblright{}.

\begin{grammarlens}[title={What you've built}]
One more word for this chapter.
\end{grammarlens}

\subsection*{Your turn}
\begin{itemize}
\raggedright
  \item \emph{Say it:} \emph{ópyt}
  \item \emph{Say it:} \emph{ópyt}, once more
  \item \emph{Say it:} say \emph{ópyt}
  \item \emph{Recall:} say the Russian for a project, then the Russian for a task; a problem, then say \emph{ópyt} again
\end{itemize}

\begin{tcolorbox}[breakable,colback=teal!4,colframe=teal!35!black,title={Before you move on}]
What does \ru{опыт} mean? (\textquotedblleft{}Experience; an experiment\textquotedblright{}.) Say it once more.
\end{tcolorbox}
\section[(dialogue)]{Words from chapters 250-253 — a first pass}

\label{lesson:RU-R257-first-pass-words-from-chapters-250-253}



\begin{quote}
Say the words for a task; a problem and experience; an experiment.
\end{quote}

\subsection*{Your turn}
Say these aloud:

\begin{itemize}
\raggedright
  \item a watermelon, a plum, pasta, a cake, a sweet
  \item a bull, a cockerel, a goose, a pigeon, a sparrow
  \item hockey, tennis, basketball, running, swimming
  \item a fairy tale, a story; history, chess, transport, a trolleybus
\end{itemize}

\begin{tcolorbox}[breakable,colback=teal!4,colframe=teal!35!black,title={Before you move on}]
A watermelon? (\textbf{\ru{арбуз}}.) A plum? (\textbf{\ru{слива}}.) Pasta? (\textbf{\ru{макароны}}.) A cake? (\textbf{\ru{торт}}.) A sweet? (\textbf{\ru{конфета}}.) A bull? (\textbf{\ru{бык}}.) A cockerel? (\textbf{\ru{петух}}.) A goose? (\textbf{\ru{гусь}}.) A pigeon? (\textbf{\ru{голубь}}.) A sparrow? (\textbf{\ru{воробей}}.) Hockey? (\textbf{\ru{хоккей}}.) Tennis? (\textbf{\ru{теннис}}.) Basketball? (\textbf{\ru{баскетбол}}.) Running? (\textbf{\ru{бег}}.) Swimming? (\textbf{\ru{плавание}}.) A fairy tale? (\textbf{\ru{сказка}}.) A story; history? (\textbf{\ru{история}}.) Chess? (\textbf{\ru{шахматы}}.) Transport? (\textbf{\ru{транспорт}}.) A trolleybus? (\textbf{\ru{троллейбус}}.)
\end{tcolorbox}
\section[(dialogue)]{Words from chapters 254-257 — a second pass}

\label{lesson:RU-R257-second-pass-words-from-chapters-254-257}



\begin{quote}
Say the words for a girl and a lad.
\end{quote}

\subsection*{Your turn}
Say these aloud:

\begin{itemize}
\raggedright
  \item a girl, a lad, a young woman, parents, a nephew
  \item a surname, an age, a number; a hotel room, a message, a conversation
  \item a professor, a journalist, an artist, a musician, an actor
  \item a company, a firm, a project, a task; a problem, experience; an experiment
\end{itemize}

\begin{tcolorbox}[breakable,colback=teal!4,colframe=teal!35!black,title={Before you move on}]
A girl? (\textbf{\ru{девочка}}.) A lad? (\textbf{\ru{парень}}.) A young woman? (\textbf{\ru{девушка}}.) Parents? (\textbf{\ru{родители}}.) A nephew? (\textbf{\ru{племянник}}.) A surname? (\textbf{\ru{фамилия}}.) An age? (\textbf{\ru{возраст}}.) A number; a hotel room? (\textbf{\ru{номер}}.) A message? (\textbf{\ru{сообщение}}.) A conversation? (\textbf{\ru{разговор}}.) A professor? (\textbf{\ru{профессор}}.) A journalist? (\textbf{\ru{журналист}}.) An artist? (\textbf{\ru{художник}}.) A musician? (\textbf{\ru{музыкант}}.) An actor? (\textbf{\ru{актёр}}.) A company? (\textbf{\ru{компания}}.) A firm? (\textbf{\ru{фирма}}.) A project? (\textbf{\ru{проект}}.) A task; a problem? (\textbf{\ru{задача}}.) Experience; an experiment? (\textbf{\ru{опыт}}.)
\end{tcolorbox}
